\documentclass[aspectratio=169]{beamer}
\usetheme{Madrid}
\usecolortheme{default}
\usepackage[utf8]{inputenc}
\usepackage{amsmath}
\usepackage{amsfonts}
\usepackage{amssymb}
\usepackage{enumerate}

\setbeamertemplate{navigation symbols}{}
\setbeamerfont{block body}{size=\small}
\setbeamerfont{alertblock body}{size=\small}

\title{CRISC Domains C and D -- 100 Real-Life Questions}
\subtitle{Risk Response, Controls, Monitoring, IT and Security -- Practical Scenarios}
\author{CRISC Mastery Series}
\date{\today}

% ------------------------------------------------------------
% \hq{number}{topic}{question}{options}{letter}{explanation}{takeaway}
% Creates one question frame and one explanation frame.
% ------------------------------------------------------------
\newcommand{\hq}[7]{%
	\begin{frame}{Q#1: #2}
		\begin{block}{Question}
			#3
		\end{block}
		\begin{itemize}
			#4
		\end{itemize}
	\end{frame}
	\begin{frame}{Q#1: Explanation}
		\begin{block}{Correct Answer: #5}
			\begin{itemize}
				#6
			\end{itemize}
		\end{block}
		\begin{alertblock}{Key Takeaway}
			#7
		\end{alertblock}
	\end{frame}
}

\begin{document}
	
	\begin{frame}
		\titlepage
	\end{frame}
	
	\begin{frame}{Instructions}
		\begin{alertblock}{How to Use This Deck}
			Think like a risk practitioner in a real organization.\\
			Focus on business context, ownership, and practical decision-making.\\
			The answer and a key takeaway are on the slide following each question.
		\end{alertblock}
	\end{frame}
	
	\begin{frame}{Outline}
		\tableofcontents
	\end{frame}
	
	% ============================================================
	\section{Risk Response and Ownership}
	% ============================================================
	
	\hq{1}{Choosing a Response}
	{A moderate risk is found on a low-value internal wiki. Fixing it would cost far more than the potential loss. What is the \textbf{MOST} appropriate response?}
	{\item[\textbf{(A)}] Mitigate it anyway to be safe
		\item[\textbf{(B)}] Buy a new cyber insurance policy just for the wiki
		\item[\textbf{(C)}] Accept the risk, with documented approval from the risk owner
		\item[\textbf{(D)}] Shut down every internal wiki in the company}
	{C}
	{\item A response should be cost-effective compared with the exposure.
		\item Acceptance is valid when the risk is within appetite and formally documented.
		\item Spending more than the potential loss is poor risk management.}
	{Never spend more on a control than the risk it addresses.}
	
	\hq{2}{Who Accepts the Risk?}
	{The security team proposes accepting a high risk that affects the sales order process. Who should \textbf{formally} approve the acceptance?}
	{\item[\textbf{(A)}] The Chief Information Security Officer
		\item[\textbf{(B)}] The internal audit manager
		\item[\textbf{(C)}] The business owner of the sales order process
		\item[\textbf{(D)}] The IT service desk manager}
	{C}
	{\item The risk owner is accountable for the business process and its outcomes.
		\item Security advises and recommends; it does not own the business risk.
		\item Audit is independent and must not accept risk on behalf of management.}
	{Risk is accepted by the process owner, not by the team that found it.}
	
	\hq{3}{Residual Risk Above Appetite}
	{After controls are applied, residual risk still exceeds the approved risk appetite. What should the risk practitioner do \textbf{NEXT}?}
	{\item[\textbf{(A)}] Lower the risk appetite so the risk fits inside it
		\item[\textbf{(B)}] Escalate to the risk owner and senior management for a decision on further response or a formal exception
		\item[\textbf{(C)}] Close the risk because controls were implemented
		\item[\textbf{(D)}] Re-rate the residual risk lower so it looks acceptable}
	{B}
	{\item Appetite is set by leadership, not adjusted to fit one risk.
		\item Residual risk beyond appetite needs a management decision.
		\item Changing ratings to hide a gap defeats the purpose of risk management.}
	{If residual risk is outside appetite, the decision goes up, not around.}
	
	\hq{4}{Cyber Insurance}
	{A company buys cyber insurance to cover ransomware losses. Which statement about the remaining risk is \textbf{MOST} accurate?}
	{\item[\textbf{(A)}] All ransomware risk has been eliminated
		\item[\textbf{(B)}] Accountability, reputational damage, and operational disruption remain with the organization
		\item[\textbf{(C)}] Accountability has moved to the insurer
		\item[\textbf{(D)}] The likelihood of an attack has dropped significantly}
	{B}
	{\item Insurance transfers part of the financial impact only.
		\item It does not stop attacks or reduce their likelihood.
		\item Customers, regulators, and the board still hold the organization accountable.}
	{You can transfer some of the cost, never the accountability.}
	
	\hq{5}{What Is Avoidance?}
	{A company wants to launch a product in a country whose data laws conflict with its own obligations and cannot be satisfied. Which action is risk \textbf{avoidance}?}
	{\item[\textbf{(A)}] Buy insurance for regulatory fines
		\item[\textbf{(B)}] Add stronger encryption and launch anyway
		\item[\textbf{(C)}] Decide not to launch the product in that country
		\item[\textbf{(D)}] Launch and monitor for complaints}
	{C}
	{\item Avoidance removes the activity that creates the risk.
		\item Insurance is transfer, encryption is mitigation, and monitoring is acceptance.
		\item Avoidance is appropriate when no other response brings risk within appetite.}
	{Avoidance means not doing the risky thing.}
	
	\hq{6}{Selecting a Response}
	{Several response options exist for the same risk. Which factor is \textbf{MOST} important when choosing among them?}
	{\item[\textbf{(A)}] Which option the vendor recommends
		\item[\textbf{(B)}] Which option is the most technically advanced
		\item[\textbf{(C)}] The cost-benefit of each option and how well it brings risk within appetite
		\item[\textbf{(D)}] Which option can be deployed fastest}
	{C}
	{\item The goal is to bring risk within appetite at an acceptable cost.
		\item Technology or speed alone does not justify a response.
		\item Vendors have a sales interest and are not neutral.}
	{Best response = within appetite at reasonable cost.}
	
	\hq{7}{Slipping Action Plans}
	{A mitigation plan was approved six months ago and its target dates keep slipping. What should the risk practitioner do?}
	{\item[\textbf{(A)}] Report the delay to the risk owner and reassess whether the current exposure is still acceptable
		\item[\textbf{(B)}] Remove the item from the risk register
		\item[\textbf{(C)}] Extend the dates quietly without telling anyone
		\item[\textbf{(D)}] Escalate to the board without talking to the owner}
	{A}
	{\item The risk exists until it is actually reduced.
		\item The owner must know the exposure and decide whether to keep waiting.
		\item Silent extensions hide ongoing exposure.}
	{A late action plan means the risk is still open.}
	
	\hq{8}{Policy Exceptions}
	{A business unit asks for a temporary exception to the password policy for a legacy application. What is the \textbf{BEST} approach?}
	{\item[\textbf{(A)}] Deny all exceptions without discussion
		\item[\textbf{(B)}] Assess the risk, require compensating controls, set an expiry date, and obtain owner approval
		\item[\textbf{(C)}] Grant it verbally to keep the business moving
		\item[\textbf{(D)}] Rewrite the policy to match the legacy application}
	{B}
	{\item Exceptions are normal but must be risk-assessed and time-bound.
		\item Compensating controls reduce exposure while the gap exists.
		\item Undocumented exceptions are unmanaged risk.}
	{An exception without an expiry date is a permanent gap.}
	
	\hq{9}{Legacy System Without MFA}
	{A critical legacy system cannot support multi-factor authentication. What should the risk practitioner recommend?}
	{\item[\textbf{(A)}] Implement compensating controls such as segmentation, restricted access, and enhanced monitoring, and document the exception
		\item[\textbf{(B)}] Accept the risk and take no action
		\item[\textbf{(C)}] Replace the system immediately regardless of cost
		\item[\textbf{(D)}] Remove the system from the asset inventory}
	{A}
	{\item Compensating controls reduce risk when the preferred control is not possible.
		\item Immediate replacement may not be cost-justified or feasible.
		\item Hiding the asset makes the risk invisible, not smaller.}
	{When the ideal control is impossible, compensate and document.}
	
	\hq{10}{Secondary Risk}
	{Moving backups to a cloud provider reduces the risk of local data loss but creates new exposure of the backup data. What is this new exposure called?}
	{\item[\textbf{(A)}] Inherent risk
		\item[\textbf{(B)}] Residual risk of the original scenario
		\item[\textbf{(C)}] Risk capacity
		\item[\textbf{(D)}] Secondary risk, which must be assessed as part of the response}
	{D}
	{\item A response can introduce new risks of its own.
		\item These must be evaluated before the response is approved.
		\item Ignoring them can leave the organization worse off.}
	{Every response can create a new risk; check for it.}
	
	\hq{11}{How Much Risk Remains?}
	{Management asks how much risk will remain once all planned controls are fully implemented. Which measure answers this?}
	{\item[\textbf{(A)}] Inherent risk
		\item[\textbf{(B)}] Residual risk
		\item[\textbf{(C)}] Risk capacity
		\item[\textbf{(D)}] Risk appetite}
	{B}
	{\item Residual risk is what is left after controls are applied.
		\item Inherent risk is the level before any controls.
		\item Appetite and capacity describe how much risk the organization is willing or able to take.}
	{Inherent minus controls equals residual.}
	
	\hq{12}{Risk With No Owner}
	{A risk has been in the register for nine months with no assigned owner. What should the risk practitioner do \textbf{FIRST}?}
	{\item[\textbf{(A)}] Work with senior management to assign an owner who is accountable for the affected business process
		\item[\textbf{(B)}] Assign it to the risk practitioner to speed things up
		\item[\textbf{(C)}] Assign it to the IT department by default
		\item[\textbf{(D)}] Close the risk because nobody is concerned}
	{A}
	{\item An owner must have authority and accountability for the process.
		\item Practitioners advise on risk; they should not own business risk.
		\item Unowned risks tend to go unmanaged.}
	{A risk without an owner is a problem, not a managed risk.}
	
	\hq{13}{Prioritizing Responses}
	{The budget is limited and many risks need action. How should responses be \textbf{prioritized}?}
	{\item[\textbf{(A)}] Cheapest fix first
		\item[\textbf{(B)}] Most technically interesting first
		\item[\textbf{(C)}] By risk level against appetite and business impact, agreed with the risk owners
		\item[\textbf{(D)}] Most recently discovered first}
	{C}
	{\item Prioritization should follow exposure and business impact.
		\item Owners must agree because they bear the consequences.
		\item Cost, novelty, or recency alone can leave critical risks open.}
	{Spend where the business exposure is highest.}
	
	\hq{14}{Validating a New Control}
	{A mitigation project has just been completed. What should the risk practitioner do \textbf{NEXT}?}
	{\item[\textbf{(A)}] Validate that the control works as intended and update the residual risk in the register
		\item[\textbf{(B)}] Close the item without any validation
		\item[\textbf{(C)}] Wait for the next annual assessment
		\item[\textbf{(D)}] Inform only the audit team}
	{A}
	{\item Implementation does not prove effectiveness.
		\item Residual risk must reflect the real control performance.
		\item Timely updates keep the register reliable for decisions.}
	{Implemented is not the same as effective.}
	
	% ============================================================
	\section{Third-Party Risk Management}
	% ============================================================
	
	\hq{15}{Onboarding a Vendor}
	{The business wants to onboard a vendor that will process customer data. What should the risk practitioner do \textbf{FIRST}?}
	{\item[\textbf{(A)}] Sign the contract to avoid delaying the project
		\item[\textbf{(B)}] Ask only for a penetration test report
		\item[\textbf{(C)}] Give the vendor access to a test environment
		\item[\textbf{(D)}] Perform due diligence and a risk assessment based on the data involved and the vendor's criticality}
	{D}
	{\item Risk depends on the data, the service, and the vendor's controls.
		\item Contracts should reflect the findings of the assessment.
		\item Access is granted only after risk is understood.}
	{Assess before you sign.}
	
	\hq{16}{Contract Provisions}
	{Which contract provision \textbf{MOST} helps the organization oversee a vendor's control environment?}
	{\item[\textbf{(A)}] A guaranteed price for three years
		\item[\textbf{(B)}] Right to audit and a requirement to provide independent assurance reports
		\item[\textbf{(C)}] An automatic renewal clause
		\item[\textbf{(D)}] A marketing statement about industry leadership}
	{B}
	{\item Oversight needs the ability to verify controls.
		\item Independent reports provide objective evidence.
		\item Commercial terms do not provide assurance about security.}
	{Trust, but verify: contract for the right to check.}
	
	\hq{17}{Fourth Parties}
	{A vendor hosts your data with a subcontractor. What is the \textbf{BEST} way to manage this exposure?}
	{\item[\textbf{(A)}] Ignore it, since the subcontractor is not your vendor
		\item[\textbf{(B)}] Ask the subcontractor to sign your employee handbook
		\item[\textbf{(C)}] Cancel the contract automatically
		\item[\textbf{(D)}] Require the vendor to disclose subcontractors and flow down your security and privacy requirements}
	{D}
	{\item You remain accountable for your data wherever it is processed.
		\item Flow-down clauses make the vendor responsible for its suppliers.
		\item Visibility into the chain is needed to assess true risk.}
	{Your data's risk includes your vendor's vendors.}
	
	\hq{18}{Qualified SOC Report}
	{A vendor's SOC 2 report contains a qualified opinion on access management. What is the \textbf{BEST} action?}
	{\item[\textbf{(A)}] Ignore it because the vendor has a SOC 2 report
		\item[\textbf{(B)}] Terminate the contract immediately without analysis
		\item[\textbf{(C)}] Accept the report as it is without review
		\item[\textbf{(D)}] Assess the impact on your own risk, discuss remediation with the vendor, and consider compensating controls}
	{D}
	{\item A qualification signals a control weakness that may affect you.
		\item The impact depends on your data and use of the service.
		\item The response should be proportionate and risk-based.}
	{Read the report, not just the cover page.}
	
	\hq{19}{Vendor Breach Notice}
	{A vendor reports a breach that may involve your customer data. What should be done \textbf{FIRST}?}
	{\item[\textbf{(A)}] Wait for the vendor to finish its investigation
		\item[\textbf{(B)}] Terminate the contract that day
		\item[\textbf{(C)}] Activate incident response, gather facts from the vendor, and assess impact against contractual and regulatory obligations
		\item[\textbf{(D)}] Notify all customers before any facts are known}
	{C}
	{\item Your own incident process should start immediately.
		\item Decisions on notification depend on verified facts and obligations.
		\item Premature action can cause harm and legal exposure.}
	{Facts first, then notification and contract decisions.}
	
	\hq{20}{Ongoing Vendor Monitoring}
	{How should the organization decide how often to reassess a critical vendor?}
	{\item[\textbf{(A)}] Only at contract signing
		\item[\textbf{(B)}] Every ten years for all vendors
		\item[\textbf{(C)}] Only after an incident occurs
		\item[\textbf{(D)}] Based on criticality and risk, with periodic reviews plus reviews triggered by significant events}
	{D}
	{\item Higher-risk vendors need more frequent oversight.
		\item Triggers include breaches, mergers, and service changes.
		\item One-time assessments become outdated quickly.}
	{Vendor risk changes, so monitoring must continue.}
	
	\hq{21}{Concentration Risk}
	{Ninety percent of critical services depend on one cloud provider. Which concern is \textbf{MOST} relevant?}
	{\item[\textbf{(A)}] Only the monthly cost of the service
		\item[\textbf{(B)}] There is no concern because the provider is large
		\item[\textbf{(C)}] The concern applies only to small vendors
		\item[\textbf{(D)}] Concentration risk, requiring an assessment of exit strategy and alternatives}
	{D}
	{\item Heavy dependence on one provider increases impact if it fails.
		\item An exit or contingency plan reduces lock-in and outage risk.
		\item Size does not guarantee resilience.}
	{Depend on one provider and you inherit its failures.}
	
	\hq{22}{Vendor Offboarding}
	{A vendor contract is ending. Which risk activity is \textbf{MOST} important at termination?}
	{\item[\textbf{(A)}] Delete the vendor from the accounting system only
		\item[\textbf{(B)}] Keep access active in case of future work
		\item[\textbf{(C)}] Confirm return or secure destruction of data and revoke all access, with evidence
		\item[\textbf{(D)}] Send a thank-you letter}
	{C}
	{\item Data and access can remain exposed after the relationship ends.
		\item Evidence of destruction supports compliance.
		\item Access removal prevents misuse of dormant accounts.}
	{The risk does not end when the contract does.}
	
	% ============================================================
	\section{Control Design, Implementation and Testing}
	% ============================================================
	
	\hq{23}{Control Types}
	{A badge reader denies building entry to anyone without authorization. What type of control is this?}
	{\item[\textbf{(A)}] Detective
		\item[\textbf{(B)}] Corrective
		\item[\textbf{(C)}] Preventive
		\item[\textbf{(D)}] Directive}
	{C}
	{\item Preventive controls stop an event before it occurs.
		\item Detective controls identify events after or while they happen.
		\item Corrective controls restore normal operation after an event.}
	{Preventive stops it, detective finds it, corrective fixes it.}
	
	\hq{24}{Reviewing Privileged Logs}
	{A manager reviews privileged access logs every day for unusual activity. This is \textbf{BEST} described as which type of control?}
	{\item[\textbf{(A)}] Detective
		\item[\textbf{(B)}] Preventive
		\item[\textbf{(C)}] Corrective
		\item[\textbf{(D)}] Compensating by definition}
	{A}
	{\item The review finds activity that has already occurred.
		\item It does not block the action itself.
		\item Detective controls are most useful when responses are timely.}
	{If it looks back to find problems, it is detective.}
	
	\hq{25}{Defense in Depth}
	{Why does the organization use several layers of security controls instead of one very strong control?}
	{\item[\textbf{(A)}] No single control is perfect, so layers reduce the chance that one failure leads to compromise
		\item[\textbf{(B)}] Auditors require at least three controls per risk
		\item[\textbf{(C)}] Layers always cost less than one strong control
		\item[\textbf{(D)}] Layers remove the need for monitoring}
	{A}
	{\item Controls can fail or be bypassed.
		\item Independent layers reduce the chance of total failure.
		\item Layers still need monitoring and maintenance.}
	{Assume any single control can fail.}
	
	\hq{26}{Shiny Tool Request}
	{Management wants to buy a popular security tool after reading a magazine review. What should the risk practitioner do \textbf{FIRST}?}
	{\item[\textbf{(A)}] Approve the purchase to show support
		\item[\textbf{(B)}] Compare only the vendor's feature lists
		\item[\textbf{(C)}] Ask a competitor for a demonstration
		\item[\textbf{(D)}] Confirm which identified risk the tool would reduce and whether the cost is justified}
	{D}
	{\item Controls should be tied to identified risks.
		\item Cost-benefit analysis validates the investment.
		\item Features do not equal risk reduction.}
	{Start with the risk, not the product.}
	
	\hq{27}{Admin Access Request}
	{A developer asks for permanent administrator rights on the production environment. What is the \textbf{BEST} response?}
	{\item[\textbf{(A)}] Grant full permanent rights to avoid delays
		\item[\textbf{(B)}] Refuse and give no alternative
		\item[\textbf{(C)}] Share an existing admin password
		\item[\textbf{(D)}] Grant only the access needed, time-limited and approved, with activity logged}
	{D}
	{\item Least privilege limits the impact of mistakes and misuse.
		\item Time-bound access reduces standing exposure.
		\item Shared credentials destroy accountability.}
	{Give the minimum access, for the minimum time.}
	
	\hq{28}{Segregation of Duties}
	{One employee can create vendor records and approve payments to them. What is the \textbf{BEST} recommendation?}
	{\item[\textbf{(A)}] Trust the employee because of long service
		\item[\textbf{(B)}] Remove all payment approvals
		\item[\textbf{(C)}] Separate the duties, or add compensating controls such as independent review of changes
		\item[\textbf{(D)}] Rely on the annual audit to detect problems}
	{C}
	{\item Combining incompatible duties enables fraud and error.
		\item Separation or independent review reduces that risk.
		\item Annual audits detect issues late.}
	{No single person should control a transaction end to end.}
	
	\hq{29}{Design vs Operation}
	{A well-designed access review control is not being performed by the managers responsible. Which type of weakness is this?}
	{\item[\textbf{(A)}] A design deficiency
		\item[\textbf{(B)}] A weakness in inherent risk
		\item[\textbf{(C)}] A scope limitation
		\item[\textbf{(D)}] An operating effectiveness weakness}
	{D}
	{\item The control exists on paper but is not carried out.
		\item A design deficiency means the control could not work even if performed.
		\item The fix is accountability and monitoring, not redesign.}
	{A good design that is not performed is still a failure.}
	
	\hq{30}{Strongest Test Evidence}
	{Which approach gives the \textbf{STRONGEST} evidence that a control operated effectively over the period?}
	{\item[\textbf{(A)}] Re-performing the control on a sample of evidence from throughout the period
		\item[\textbf{(B)}] Interviewing the control owner
		\item[\textbf{(C)}] Reading the control description
		\item[\textbf{(D)}] Obtaining management's written assertion}
	{A}
	{\item Re-performance and inspection rely on actual evidence.
		\item Interviews and assertions are weaker, self-reported evidence.
		\item Samples should cover the whole review period.}
	{Evidence beats explanations.}
	
	\hq{31}{Post-Implementation Review}
	{What is the primary purpose of a post-implementation review of a new control?}
	{\item[\textbf{(A)}] Assign blame for project delays
		\item[\textbf{(B)}] Verify that the control achieves its intended risk reduction and meets requirements
		\item[\textbf{(C)}] Prepare the vendor invoice for payment
		\item[\textbf{(D)}] Update the organization chart}
	{B}
	{\item The review confirms the control delivers the expected benefits.
		\item Findings feed lessons learned and residual risk updates.
		\item Its focus is value and effectiveness, not blame.}
	{Check that the control delivered what was promised.}
	
	\hq{32}{Ownerless Control}
	{A key control fails repeatedly and nobody takes responsibility. What is the \textbf{BEST} action?}
	{\item[\textbf{(A)}] Replace the control with a policy statement
		\item[\textbf{(B)}] Remove the control to avoid false assurance
		\item[\textbf{(C)}] Have audit operate the control
		\item[\textbf{(D)}] Assign a control owner with defined responsibilities and performance measures}
	{D}
	{\item Controls need accountable owners.
		\item Defined measures show when the control is not working.
		\item Audit must stay independent of operating controls.}
	{No owner, no control.}
	
	\hq{33}{Manual Access Reviews}
	{Manual quarterly access reviews are repeatedly missed. Which improvement is \textbf{MOST} effective?}
	{\item[\textbf{(A)}] Automate scheduling and tracking, for example with an identity governance tool, and monitor completion
		\item[\textbf{(B)}] Send a stronger reminder email once a year
		\item[\textbf{(C)}] Cancel the reviews
		\item[\textbf{(D)}] Ask audit to complete them}
	{A}
	{\item Automation improves consistency and evidence.
		\item Monitoring completion shows whether reviews are done.
		\item Cancelling or outsourcing to audit removes the control or independence.}
	{Automate repeatable controls where practical.}
	
	\hq{34}{Secure Baselines}
	{New servers are built inconsistently, creating random security gaps. Which control is \textbf{BEST}?}
	{\item[\textbf{(A)}] Standard secure configuration baselines with automated compliance checking
		\item[\textbf{(B)}] A longer written policy
		\item[\textbf{(C)}] More frequent vulnerability meetings
		\item[\textbf{(D)}] Let each administrator decide the settings}
	{A}
	{\item Baselines give a repeatable secure starting state.
		\item Automated checks detect drift.
		\item Policy alone does not enforce consistent builds.}
	{Standardize the build; do not rely on memory.}
	
	\hq{35}{Reorganization and Controls}
	{After a reorganization, several controls no longer match how work is performed. What should the risk practitioner do?}
	{\item[\textbf{(A)}] Leave the controls as they are
		\item[\textbf{(B)}] Reassess control design against current processes and risks, and update where needed
		\item[\textbf{(C)}] Remove all controls tied to the old structure
		\item[\textbf{(D)}] Wait for audit to find the gaps}
	{B}
	{\item Controls must fit the current process.
		\item Business change often creates gaps and overlaps.
		\item Proactive review avoids audit surprises.}
	{Change the business, then recheck the controls.}
	
	\hq{36}{Policy Not Followed}
	{Policy requires laptop encryption, but a review finds many unencrypted laptops. How should this be treated?}
	{\item[\textbf{(A)}] As a documentation error only
		\item[\textbf{(B)}] As an acceptable condition because the policy is old
		\item[\textbf{(C)}] As a control gap that needs risk assessment and a response plan
		\item[\textbf{(D)}] As a matter for HR only}
	{C}
	{\item Non-compliance with policy creates exposure.
		\item Impact depends on the data stored and the likelihood of loss.
		\item The response should be owned and tracked.}
	{A policy is only as good as its enforcement.}
	
	% ============================================================
	\section{Monitoring, Logging and Reporting}
	% ============================================================
	
	\hq{37}{Why a SIEM?}
	{What is the primary value of a SIEM to the risk program?}
	{\item[\textbf{(A)}] It aggregates and correlates logs from many sources to detect anomalies and support investigations
		\item[\textbf{(B)}] It replaces the need for firewalls
		\item[\textbf{(C)}] It prevents all attacks automatically
		\item[\textbf{(D)}] It stores only employee emails}
	{A}
	{\item Correlation reveals patterns that single logs miss.
		\item It supports incident response and compliance evidence.
		\item It is a detective and analysis tool, not a complete defense.}
	{A SIEM finds the pattern, not the whole cure.}
	
	\hq{38}{Protecting Logs}
	{Which measure \textbf{BEST} protects logs from tampering by an intruder?}
	{\item[\textbf{(A)}] Store logs only on the server that generated them
		\item[\textbf{(B)}] Allow all administrators to edit logs
		\item[\textbf{(C)}] Delete logs weekly
		\item[\textbf{(D)}] Forward logs to centralized, write-protected storage with restricted access}
	{D}
	{\item Attackers often try to alter local logs to hide activity.
		\item Central, protected storage preserves evidence.
		\item Limiting access supports accountability.}
	{Logs are evidence; protect them like evidence.}
	
	\hq{39}{Alert Fatigue}
	{The monitoring team ignores most alerts because of high false positives. What is the \textbf{BEST} improvement?}
	{\item[\textbf{(A)}] Turn off alerting for noisy systems
		\item[\textbf{(B)}] Tune detection rules and prioritize alerts by risk
		\item[\textbf{(C)}] Hire more staff only to read every alert
		\item[\textbf{(D)}] Lower all alert thresholds}
	{B}
	{\item Tuning improves signal quality.
		\item Risk-based prioritization focuses attention on what matters.
		\item Disabling or flooding alerts increases the chance of missing real events.}
	{Reduce the noise before you add people.}
	
	\hq{40}{Log Retention}
	{How should the retention period for logs be determined?}
	{\item[\textbf{(A)}] By available disk space only
		\item[\textbf{(B)}] By the administrator's preference
		\item[\textbf{(C)}] By the shortest period that auditors have ever requested
		\item[\textbf{(D)}] Based on legal, regulatory, investigative, and business requirements}
	{D}
	{\item Different laws and contracts require different periods.
		\item Investigations may need older data.
		\item Retention should be defined in policy.}
	{Retention is a policy decision, not a storage decision.}
	
	\hq{41}{Setting KRI Thresholds}
	{How should KRI thresholds be \textbf{BEST} set?}
	{\item[\textbf{(A)}] At whatever level produces few alerts
		\item[\textbf{(B)}] By copying another company's values
		\item[\textbf{(C)}] At the level the technology team prefers
		\item[\textbf{(D)}] Aligned with the risk appetite and tolerance approved by management}
	{D}
	{\item Thresholds define when risk becomes unacceptable.
		\item They should reflect decisions on appetite and tolerance.
		\item Arbitrary thresholds create false comfort or constant alarms.}
	{A threshold is appetite translated into numbers.}
	
	\hq{42}{Too Many KRIs}
	{A dashboard shows 80 KRIs and executives cannot tell which matter. What is the \textbf{BEST} action?}
	{\item[\textbf{(A)}] Add more indicators for completeness
		\item[\textbf{(B)}] Hide the dashboard from executives
		\item[\textbf{(C)}] Reduce to a manageable set linked to key risks and the decisions they support
		\item[\textbf{(D)}] Replace indicators with a single color}
	{C}
	{\item Indicators should drive decisions.
		\item A focused set is easier to monitor and act on.
		\item Removing visibility or detail does not help management.}
	{Measure what you will act on.}
	
	\hq{43}{A Leading Indicator}
	{Which metric is the \textbf{BEST} leading indicator for phishing risk?}
	{\item[\textbf{(A)}] The number of past breaches last year
		\item[\textbf{(B)}] A rising click rate in simulated phishing campaigns
		\item[\textbf{(C)}] Total annual loss from fraud
		\item[\textbf{(D)}] Average time to close helpdesk tickets}
	{B}
	{\item Leading indicators signal increasing exposure before incidents occur.
		\item Past breaches and losses are lagging measures.
		\item Ticket time is unrelated to phishing susceptibility.}
	{Leading indicators warn; lagging indicators report.}
	
	\hq{44}{Control Self-Assessment}
	{What is a key benefit of control self-assessment (CSA)?}
	{\item[\textbf{(A)}] It replaces independent audit completely
		\item[\textbf{(B)}] It involves process owners in evaluating their own controls and builds ownership
		\item[\textbf{(C)}] It guarantees an unqualified opinion
		\item[\textbf{(D)}] It eliminates the need for testing}
	{B}
	{\item CSA increases awareness and accountability in the first line.
		\item Results should still be validated independently.
		\item It supplements rather than replaces audit.}
	{CSA builds ownership; it does not replace independence.}
	
	\hq{45}{Continuous Monitoring}
	{Why is continuous monitoring preferred over only periodic assessments for critical risks?}
	{\item[\textbf{(A)}] It is always cheaper than assessments
		\item[\textbf{(B)}] It detects changes in risk and control performance sooner, supporting timely response
		\item[\textbf{(C)}] It removes the need for risk owners
		\item[\textbf{(D)}] It is required for all risks regardless of importance}
	{B}
	{\item Risk can change between assessment cycles.
		\item Timely information supports faster decisions.
		\item Intensity should match criticality.}
	{Fast-moving risks need fast-moving information.}
	
	\hq{46}{Tailoring Reports}
	{A risk practitioner sends the same detailed technical report to engineers and the board. What is the \textbf{BEST} improvement?}
	{\item[\textbf{(A)}] Send the board more technical detail
		\item[\textbf{(B)}] Send engineers only a one-line summary
		\item[\textbf{(C)}] Stop reporting to the board
		\item[\textbf{(D)}] Tailor the content and level of detail to each audience's decisions and needs}
	{D}
	{\item Boards need risk versus appetite and decisions required.
		\item Engineers need actionable technical detail.
		\item One report cannot serve everyone.}
	{Know your audience, then write.}
	
	\hq{47}{Conflicting Metrics}
	{Risk metrics from two teams for the same control disagree. What should the practitioner do \textbf{FIRST}?}
	{\item[\textbf{(A)}] Report the more favorable figure
		\item[\textbf{(B)}] Average the two numbers
		\item[\textbf{(C)}] Validate data sources and standardize definitions before reporting
		\item[\textbf{(D)}] Remove the metric from the report}
	{C}
	{\item Reports depend on reliable data.
		\item Common definitions make results comparable.
		\item Choosing the convenient number damages credibility.}
	{Bad data makes a good report dangerous.}
	
	\hq{48}{Triggers for Updates}
	{The company announces an acquisition. What should the risk practitioner do regarding the risk register?}
	{\item[\textbf{(A)}] Wait until the next annual review
		\item[\textbf{(B)}] Do nothing because acquisitions are not IT events
		\item[\textbf{(C)}] Delete low risks to simplify the register
		\item[\textbf{(D)}] Review and update risks and ownership to reflect the new environment}
	{D}
	{\item Business changes can create new risks and change existing ones.
		\item Early updates keep the register decision-ready.
		\item Risk management is continuous, not annual.}
	{Major business change means a risk refresh.}
	
	\hq{49}{Overdue Audit Findings}
	{Remediation of a high-priority audit finding is overdue. What is the \textbf{BEST} action?}
	{\item[\textbf{(A)}] Remove the finding from the report
		\item[\textbf{(B)}] Extend the due date without review
		\item[\textbf{(C)}] Accept the risk on behalf of the owner
		\item[\textbf{(D)}] Escalate per governance procedures and confirm with the owner whether the exposure is still acceptable}
	{D}
	{\item Overdue items reflect continuing exposure.
		\item Escalation ensures management is aware.
		\item Only the owner can accept risk.}
	{Overdue findings are risks that someone must decide on.}
	
	\hq{50}{Measuring Monitoring}
	{Which metric \textbf{BEST} indicates how effective an organization's monitoring capability is?}
	{\item[\textbf{(A)}] Number of monitoring tools purchased
		\item[\textbf{(B)}] Mean time to detect and respond to security incidents
		\item[\textbf{(C)}] Number of log sources in the data center
		\item[\textbf{(D)}] Total size of the log archive}
	{B}
	{\item Detection and response times show real effectiveness.
		\item Tools, log volume, and storage measure activity, not results.
		\item Lower times indicate that issues are found and handled sooner.}
	{Measure outcomes, not inventory.}
	
	% ============================================================
	\section{Architecture, Networks and Cloud}
	% ============================================================
	
	\hq{51}{Enterprise Architecture}
	{How does an enterprise architecture view help the risk practitioner?}
	{\item[\textbf{(A)}] It replaces the risk register
		\item[\textbf{(B)}] It shows systems, data flows, and dependencies so risk concentrations can be identified
		\item[\textbf{(C)}] It guarantees compliance
		\item[\textbf{(D)}] It is only useful for network engineers}
	{B}
	{\item Architecture documents how components connect and depend on each other.
		\item Dependencies reveal single points of failure.
		\item It supports impact analysis for changes.}
	{You cannot manage risk in a system you cannot map.}
	
	\hq{52}{Capability Maturity}
	{A maturity assessment rates a key process at the lowest level, often called initial. What does this mean?}
	{\item[\textbf{(A)}] The process is fully optimized
		\item[\textbf{(B)}] The process is quantitatively managed
		\item[\textbf{(C)}] The process has been independently certified
		\item[\textbf{(D)}] The process is ad hoc, unpredictable, and dependent on individuals}
	{D}
	{\item Low maturity means results are inconsistent.
		\item Higher levels bring documentation, measurement, and improvement.
		\item Maturity helps prioritize process investment.}
	{Low maturity means heroics, not repeatability.}
	
	\hq{53}{Flat Network}
	{Point-of-sale terminals share the same network as general office computers. What is the \textbf{BEST} recommendation?}
	{\item[\textbf{(A)}] Install more antivirus on office computers only
		\item[\textbf{(B)}] Segment the network so a compromise in one zone cannot easily spread to others
		\item[\textbf{(C)}] Rename the network
		\item[\textbf{(D)}] Increase the Wi-Fi password length}
	{B}
	{\item Segmentation limits lateral movement.
		\item Sensitive systems should sit in controlled zones.
		\item Endpoint tools alone do not contain network-wide compromise.}
	{Separate what is critical from what is convenient.}
	
	\hq{54}{Firewall Rule Sprawl}
	{A firewall contains hundreds of old rules that no one can explain. What is the \textbf{BEST} action?}
	{\item[\textbf{(A)}] Leave the rules, because deleting is risky
		\item[\textbf{(B)}] Add a rule allowing all traffic temporarily
		\item[\textbf{(C)}] Perform a regular rule review and remove unneeded rules through change management
		\item[\textbf{(D)}] Disable logging to reduce clutter}
	{C}
	{\item Unused rules widen the attack surface.
		\item Changes must be controlled and tested.
		\item Logging is needed to understand what is used.}
	{Review the rule base like you review user access.}
	
	\hq{55}{IDS vs IPS}
	{Which statement correctly distinguishes an IDS from an IPS?}
	{\item[\textbf{(A)}] An IDS detects and alerts, while an IPS can also block suspicious traffic
		\item[\textbf{(B)}] An IDS blocks traffic, while an IPS only logs
		\item[\textbf{(C)}] They are identical devices
		\item[\textbf{(D)}] An IPS is used only for wireless networks}
	{A}
	{\item Detection systems generate alerts for analysis.
		\item Prevention systems are placed inline and can act automatically.
		\item Both need tuning to limit false positives.}
	{IDS watches; IPS watches and acts.}
	
	\hq{56}{Fake Website}
	{Users type the correct web address but land on a fraudulent site. Which attack is \textbf{MOST} likely?}
	{\item[\textbf{(A)}] Brute-force password guessing
		\item[\textbf{(B)}] DNS poisoning or spoofing
		\item[\textbf{(C)}] Denial of service
		\item[\textbf{(D)}] SQL injection against the user's laptop}
	{B}
	{\item Corrupted DNS answers redirect users to attacker sites.
		\item Controls include secure resolvers and DNS security extensions.
		\item Monitoring unusual DNS records helps detection.}
	{If the address is right but the site is wrong, suspect DNS.}
	
	\hq{57}{Guest Wi-Fi}
	{Guests connect to a Wi-Fi network that also reaches internal servers. What is the \textbf{BEST} fix?}
	{\item[\textbf{(A)}] Place guest access on a separate, isolated network with strong encryption and authentication
		\item[\textbf{(B)}] Hide the network name
		\item[\textbf{(C)}] Ask guests to be careful
		\item[\textbf{(D)}] Move the servers to a new building}
	{A}
	{\item Isolation prevents guests from reaching internal systems.
		\item Encryption protects traffic.
		\item Hiding a network name is not a security control.}
	{Guests get internet, not the internal network.}
	
	\hq{58}{Remote Access}
	{Employees work from public networks. Which control \textbf{BEST} protects their connection to company systems?}
	{\item[\textbf{(A)}] A shared password sent by email
		\item[\textbf{(B)}] Disabling endpoint updates
		\item[\textbf{(C)}] A VPN with multi-factor authentication
		\item[\textbf{(D)}] Allowing direct internet access to servers}
	{C}
	{\item A VPN encrypts traffic over untrusted networks.
		\item MFA reduces the impact of stolen credentials.
		\item Direct exposure increases attack surface.}
	{Treat every public network as hostile.}
	
	\hq{59}{Shared Responsibility: IaaS}
	{In an infrastructure-as-a-service model, who is responsible for patching the guest operating system?}
	{\item[\textbf{(A)}] The cloud provider
		\item[\textbf{(B)}] The customer
		\item[\textbf{(C)}] The internet service provider
		\item[\textbf{(D)}] The hardware manufacturer}
	{B}
	{\item The provider secures the underlying infrastructure.
		\item The customer manages what it deploys on top, including operating systems.
		\item Responsibilities should be confirmed in the contract.}
	{The more control you have, the more you must secure.}
	
	\hq{60}{Shared Responsibility: SaaS}
	{Even when using a SaaS application, what does the customer \textbf{ALWAYS} remain responsible for?}
	{\item[\textbf{(A)}] Data classification, user access management, and configuration choices
		\item[\textbf{(B)}] Physical data center security
		\item[\textbf{(C)}] Hypervisor patching
		\item[\textbf{(D)}] Provider network hardware}
	{A}
	{\item SaaS shifts infrastructure duties to the provider.
		\item The customer still decides who can access which data.
		\item Accountability for the data remains with the customer.}
	{You can outsource the service, not the data accountability.}
	
	\hq{61}{Cloud Exposure Root Cause}
	{What is the \textbf{MOST} common root cause of data exposures in cloud environments?}
	{\item[\textbf{(A)}] Physical theft of provider hardware
		\item[\textbf{(B)}] Rare failures in encryption mathematics
		\item[\textbf{(C)}] Lack of electricity
		\item[\textbf{(D)}] Misconfiguration of storage, identity, or network settings}
	{D}
	{\item Default or careless settings can make data public.
		\item Continuous configuration monitoring catches drift.
		\item Least privilege limits damage when mistakes happen.}
	{Most cloud breaches start with a setting, not a genius.}
	
	\hq{62}{Multi-Tenancy}
	{Which risk is \textbf{MOST} specific to multi-tenant public cloud environments?}
	{\item[\textbf{(A)}] Employees working from home
		\item[\textbf{(B)}] Printer supply shortages
		\item[\textbf{(C)}] Failure of logical separation between customers' data and workloads
		\item[\textbf{(D)}] Lack of paper records}
	{C}
	{\item Tenants share underlying infrastructure.
		\item Isolation depends on the provider's design and configuration.
		\item Encryption and assurance reports help manage this risk.}
	{Shared infrastructure demands strong isolation.}
	
	\hq{63}{Cloud Logging}
	{The organization wants visibility into activity in a cloud service. What is the \textbf{BEST} approach?}
	{\item[\textbf{(A)}] Assume the provider will alert the customer
		\item[\textbf{(B)}] Disable logs to reduce cost
		\item[\textbf{(C)}] Review logs once at year end
		\item[\textbf{(D)}] Enable cloud activity logs and integrate them into central monitoring, agreeing access in the contract}
	{D}
	{\item The customer is responsible for monitoring its own use.
		\item Central integration enables correlation with other sources.
		\item Retention and access should be contractually defined.}
	{If you cannot see it, you cannot monitor it.}
	
	\hq{64}{Zero Trust}
	{Which statement \textbf{BEST} describes the zero trust principle?}
	{\item[\textbf{(A)}] Trust everyone inside the office network
		\item[\textbf{(B)}] Trust devices that were approved last year
		\item[\textbf{(C)}] Trust any user with a company email address
		\item[\textbf{(D)}] Never trust implicitly based on network location; always verify identity and context}
	{D}
	{\item Location alone is not proof of trustworthiness.
		\item Access decisions use identity, device health, and context.
		\item It limits the damage when one account is compromised.}
	{Verify every request, every time.}
	
	% ============================================================
	\section{Resilience, Data and System Development}
	% ============================================================
	
	\hq{65}{Start With a BIA}
	{Before writing a business continuity plan, what should the organization do \textbf{FIRST}?}
	{\item[\textbf{(A)}] Buy a secondary data center
		\item[\textbf{(B)}] Write recovery scripts for all systems
		\item[\textbf{(C)}] Perform a business impact analysis to identify critical processes and recovery needs
		\item[\textbf{(D)}] Hire a continuity consultant before understanding processes}
	{C}
	{\item The BIA identifies what matters most and how long it can be down.
		\item Recovery strategies depend on the results.
		\item Otherwise investment may not match business need.}
	{Know what is critical before you plan to protect it.}
	
	\hq{66}{Recovery Point Objective}
	{Management says the organization can tolerate losing at most 15 minutes of transactions. Which objective does this define?}
	{\item[\textbf{(A)}] Recovery time objective (RTO)
		\item[\textbf{(B)}] Recovery point objective (RPO)
		\item[\textbf{(C)}] Maximum tolerable outage
		\item[\textbf{(D)}] Service level agreement}
	{B}
	{\item RPO describes acceptable data loss measured in time.
		\item It drives backup and replication frequency.
		\item RTO describes how quickly service must be restored.}
	{RPO is about data, RTO is about time to recover.}
	
	\hq{67}{Recovery Time Objective}
	{The order system must be back in service within four hours of an outage. Which objective does this describe?}
	{\item[\textbf{(A)}] Recovery time objective (RTO)
		\item[\textbf{(B)}] Recovery point objective (RPO)
		\item[\textbf{(C)}] Risk appetite
		\item[\textbf{(D)}] Key risk indicator}
	{A}
	{\item RTO defines how quickly a service must be restored.
		\item It drives recovery strategy and cost.
		\item RPO covers the amount of data loss allowed.}
	{RTO is the clock, RPO is the checkpoint.}
	
	\hq{68}{Tabletop Only}
	{The disaster recovery plan has been tested only by tabletop discussion for three years. What is the \textbf{MOST} significant concern?}
	{\item[\textbf{(A)}] The plan may not work technically, because actual recovery has never been demonstrated
		\item[\textbf{(B)}] Tabletop tests are prohibited
		\item[\textbf{(C)}] The plan is too short
		\item[\textbf{(D)}] The team talks too much}
	{A}
	{\item Discussion tests logic, not technical recovery.
		\item Functional tests prove that systems can actually be restored.
		\item Test results should drive plan updates.}
	{A plan never exercised is a hope, not a capability.}
	
	\hq{69}{Untested Backups}
	{Backups run nightly, but restores have never been tried. What should the risk practitioner recommend?}
	{\item[\textbf{(A)}] Increase backup frequency only
		\item[\textbf{(B)}] Perform periodic restore tests and record the results
		\item[\textbf{(C)}] Trust the backup software reports
		\item[\textbf{(D)}] Encrypt the backup logs}
	{B}
	{\item A backup is only valuable if it can be restored.
		\item Tests reveal corruption, missing data, or procedure gaps.
		\item Success reports do not prove recoverability.}
	{A backup you cannot restore is not a backup.}
	
	\hq{70}{Alternate Site Choice}
	{The RTO for a critical service is a few minutes. Which alternate site strategy is \textbf{MOST} appropriate?}
	{\item[\textbf{(A)}] A hot site or active-active arrangement
		\item[\textbf{(B)}] A cold site
		\item[\textbf{(C)}] Manual workarounds only
		\item[\textbf{(D)}] A storage unit with spare hardware}
	{A}
	{\item Short RTOs require systems already running or ready to take over.
		\item Cold sites take much longer to become operational.
		\item Cost must be weighed against business impact.}
	{The tighter the RTO, the hotter the site.}
	
	\hq{71}{Classify First}
	{What should the organization do \textbf{FIRST} to apply appropriate protection to its data?}
	{\item[\textbf{(A)}] Classify the data according to sensitivity and business value
		\item[\textbf{(B)}] Encrypt every file with the same key
		\item[\textbf{(C)}] Move everything to the cloud
		\item[\textbf{(D)}] Delete all old records}
	{A}
	{\item Classification drives which controls are appropriate.
		\item Uniform treatment wastes money or leaves critical data under-protected.
		\item Owners should confirm classifications.}
	{You cannot protect data properly until you know what it is.}
	
	\hq{72}{Data Hoarding}
	{Customer records from over ten years ago are kept indefinitely with no business need. What is the \textbf{BEST} action?}
	{\item[\textbf{(A)}] Keep everything because storage is cheap
		\item[\textbf{(B)}] Move the data to personal drives
		\item[\textbf{(C)}] Define a retention schedule and securely dispose of data that is no longer required
		\item[\textbf{(D)}] Rename the files}
	{C}
	{\item Unneeded data increases breach impact and compliance exposure.
		\item Retention should follow legal and business needs.
		\item Disposal must be secure and documented.}
	{Data you do not have cannot be stolen.}
	
	\hq{73}{Who Classifies Data?}
	{Who is \textbf{PRIMARILY} responsible for deciding how a business dataset is classified?}
	{\item[\textbf{(A)}] The data custodian or system administrator
		\item[\textbf{(B)}] The vendor hosting the data
		\item[\textbf{(C)}] The data owner
		\item[\textbf{(D)}] The help desk}
	{C}
	{\item Owners understand business value and legal obligations.
		\item Custodians implement protection according to the classification.
		\item Responsibility should be recorded.}
	{Owners decide, custodians implement.}
	
	\hq{74}{Security Early in Development}
	{At which point in the development life cycle is it generally \textbf{MOST} cost-effective to address security?}
	{\item[\textbf{(A)}] After go-live
		\item[\textbf{(B)}] During requirements and design
		\item[\textbf{(C)}] During annual audit
		\item[\textbf{(D)}] After the first incident}
	{B}
	{\item Early fixes are cheaper than rework later.
		\item Requirements should include security and privacy needs.
		\item Late discovery can delay or stop deployment.}
	{Shift left: design it in, do not bolt it on.}
	
	\hq{75}{Fast Releases}
	{An agile team releases weekly and skips security review to keep pace. What is the \textbf{BEST} solution?}
	{\item[\textbf{(A)}] Stop all releases
		\item[\textbf{(B)}] Review security only after each year
		\item[\textbf{(C)}] Integrate automated security testing and checks into the delivery pipeline
		\item[\textbf{(D)}] Let developers self-certify without testing}
	{C}
	{\item Automation scales with speed.
		\item Early feedback helps developers fix issues quickly.
		\item Skipping checks accumulates hidden risk.}
	{Make security part of the pipeline, not a roadblock.}
	
	\hq{76}{Test Environment Data}
	{A project copies live customer data into a test environment with weaker controls. What is the \textbf{BEST} recommendation?}
	{\item[\textbf{(A)}] Allow it because it is only a test
		\item[\textbf{(B)}] Use masked or synthetic data for testing
		\item[\textbf{(C)}] Give testers production admin rights
		\item[\textbf{(D)}] Disable test logging}
	{B}
	{\item Test systems often have weaker controls.
		\item Masking preserves usefulness while reducing exposure.
		\item Privacy obligations apply to copies of data.}
	{Real data in test means real risk in test.}
	
	\hq{77}{Squeezed Project Testing}
	{A project is behind schedule, and management proposes cutting security testing. What should the risk practitioner do?}
	{\item[\textbf{(A)}] Approve quietly to support the schedule
		\item[\textbf{(B)}] Explain the risk impact to the sponsor and let the accountable owner decide, with the decision documented
		\item[\textbf{(C)}] Refuse and leave the project
		\item[\textbf{(D)}] Complete the testing secretly}
	{B}
	{\item Schedule pressure creates risk trade-offs.
		\item The business owner should decide knowingly.
		\item Documentation preserves accountability.}
	{Make the trade-off visible, then let the owner decide.}
	
	\hq{78}{System Accreditation}
	{What does formal accreditation of a system typically represent?}
	{\item[\textbf{(A)}] Management's formal acceptance of the residual risk of operating the system
		\item[\textbf{(B)}] Proof that the system has no vulnerabilities
		\item[\textbf{(C)}] A guarantee from the vendor
		\item[\textbf{(D)}] A budget approval for new hardware}
	{A}
	{\item Certification evaluates controls; accreditation is the management decision.
		\item Acceptance should be based on documented risk information.
		\item Accountability stays with management.}
	{Accreditation is a risk decision, not a guarantee.}
	
	\hq{79}{Open-Source Components}
	{Developers rely heavily on open-source libraries. Which control \textbf{BEST} manages the associated risk?}
	{\item[\textbf{(A)}] Track components in an inventory and scan them for known vulnerabilities and license issues
		\item[\textbf{(B)}] Ban all open-source software with no analysis
		\item[\textbf{(C)}] Download libraries from any source
		\item[\textbf{(D)}] Assume popular libraries are always safe}
	{A}
	{\item Third-party code can contain known vulnerabilities.
		\item An inventory enables fast response when issues appear.
		\item License terms can create legal obligations.}
	{Know what is inside your software.}
	
	\hq{80}{End-of-Life Systems}
	{A production system runs an operating system the vendor no longer supports. What is the \textbf{BEST} initial response?}
	{\item[\textbf{(A)}] Ignore it, because the system still works
		\item[\textbf{(B)}] Turn off all monitoring
		\item[\textbf{(C)}] Assess the risk, plan an upgrade, and apply compensating controls until then
		\item[\textbf{(D)}] Remove the system from inventory}
	{C}
	{\item Unsupported software no longer receives security patches.
		\item Compensating controls limit exposure during transition.
		\item A documented plan and owner reduce uncertainty.}
	{No patches means rising risk over time.}
	
	% ============================================================
	\section{Emerging Technology, Security and Privacy}
	% ============================================================
	
	\hq{81}{Bring Your Own Device}
	{Employees want to use personal phones for work email. Which approach \textbf{BEST} manages the risk?}
	{\item[\textbf{(A)}] Ban phones entirely without alternatives
		\item[\textbf{(B)}] Allow any device with no controls
		\item[\textbf{(C)}] A BYOD policy with mobile device management, data separation, and user acknowledgment
		\item[\textbf{(D)}] Ask employees to be careful}
	{C}
	{\item Policy sets expectations and responsibilities.
		\item Management tools allow remote wipe and separation of work data.
		\item Blanket bans often lead to shadow practices.}
	{Enable the business, but control the data.}
	
	\hq{82}{Default Credentials on IoT}
	{Hundreds of IoT sensors were installed with default passwords. What is the \textbf{BEST} action?}
	{\item[\textbf{(A)}] Leave them because they are small devices
		\item[\textbf{(B)}] Connect them directly to the finance network
		\item[\textbf{(C)}] Change defaults, keep an inventory, and place devices on a segmented network
		\item[\textbf{(D)}] Disable all logging}
	{C}
	{\item Default credentials are widely known.
		\item Segmentation limits what a compromised device can reach.
		\item An inventory is needed to manage updates.}
	{Small devices can still be the way in.}
	
	\hq{83}{AI Decision Model}
	{The business plans to use an AI model for credit decisions. Which concern should the risk practitioner emphasize \textbf{FIRST}?}
	{\item[\textbf{(A)}] Bias, explainability, data quality, and human oversight within an AI governance process
		\item[\textbf{(B)}] The color scheme of the user interface
		\item[\textbf{(C)}] The number of servers the model runs on
		\item[\textbf{(D)}] The vendor's brand reputation only}
	{A}
	{\item Automated decisions can affect individuals unfairly.
		\item Poor data leads to poor and biased outcomes.
		\item Governance and accountability must be defined before use.}
	{AI decisions still need a human owner.}
	
	\hq{84}{Generative AI Data Leakage}
	{Staff are pasting confidential documents into a public generative AI tool. What is the \textbf{BEST} response?}
	{\item[\textbf{(A)}] Ignore it because the tool is popular
		\item[\textbf{(B)}] Publish a policy and take no further action
		\item[\textbf{(C)}] Block the whole internet
		\item[\textbf{(D)}] Set an acceptable use policy, offer an approved tool, and apply technical controls such as data loss prevention}
	{D}
	{\item Policy gives direction; tools give enforcement.
		\item An approved alternative reduces shadow use.
		\item Overblocking harms productivity.}
	{Give people a safe way to do what they will do anyway.}
	
	\hq{85}{Blockchain Risk}
	{Which risk is \textbf{MOST} relevant when an organization records business data on a blockchain?}
	{\item[\textbf{(A)}] Records are difficult to correct once written, and key management or smart contract flaws can cause loss
		\item[\textbf{(B)}] Blockchains cannot store any data
		\item[\textbf{(C)}] Blockchains are always faster than databases
		\item[\textbf{(D)}] Blockchains require no governance}
	{A}
	{\item Immutability makes mistakes and privacy requests hard to fix.
		\item Lost or stolen keys can cause irreversible loss.
		\item Smart contract code needs review and testing.}
	{Immutable cuts both ways.}
	
	\hq{86}{Quantum Computing}
	{Sensitive data must stay confidential for 25 years. What should the risk practitioner recommend regarding quantum computing?}
	{\item[\textbf{(A)}] Ignore it until quantum computers are common
		\item[\textbf{(B)}] Remove encryption altogether
		\item[\textbf{(C)}] Use shorter passwords
		\item[\textbf{(D)}] Inventory cryptography in use and plan migration to quantum-resistant algorithms}
	{D}
	{\item Adversaries may store data now and decrypt it later.
		\item Knowing where cryptography is used makes migration possible.
		\item Long-lived data needs earlier planning.}
	{Long-lived secrets need future-proof protection.}
	
	\hq{87}{Leaver Accounts}
	{A terminated employee's account was still active two months later. What is the \textbf{BEST} long-term fix?}
	{\item[\textbf{(A)}] Remind managers by email once a year
		\item[\textbf{(B)}] Disable all accounts monthly
		\item[\textbf{(C)}] Integrate HR termination events with automated account deprovisioning and review exceptions
		\item[\textbf{(D)}] Rely on the employee to log out}
	{C}
	{\item HR events are the most reliable trigger.
		\item Automation shortens the exposure window.
		\item Periodic reviews catch what automation misses.}
	{Access should end when employment ends.}
	
	\hq{88}{What Is MFA?}
	{Which of the following is an example of multi-factor authentication?}
	{\item[\textbf{(A)}] Two different passwords
		\item[\textbf{(B)}] A password and a security question
		\item[\textbf{(C)}] A password combined with a one-time code from an authenticator app
		\item[\textbf{(D)}] A long password}
	{C}
	{\item MFA uses factors from different categories: know, have, or are.
		\item Passwords and security questions are both something you know.
		\item Different factors make a single compromise less damaging.}
	{Two of the same kind is not multi-factor.}
	
	\hq{89}{Shared Admin Account}
	{System administrators share one generic admin account. What is the \textbf{BEST} improvement?}
	{\item[\textbf{(A)}] Change the shared password every year
		\item[\textbf{(B)}] Individual accountable accounts with privileged access management and session logging
		\item[\textbf{(C)}] Write the password on a secure note
		\item[\textbf{(D)}] Increase the password length only}
	{B}
	{\item Shared accounts remove accountability.
		\item Individual accounts and logs show who did what.
		\item Privileged access tools can enforce approval and recording.}
	{If everyone is admin, no one is accountable.}
	
	\hq{90}{Bulk Encryption}
	{Which type of encryption is \textbf{MOST} suitable for encrypting large volumes of data efficiently?}
	{\item[\textbf{(A)}] Asymmetric encryption
		\item[\textbf{(B)}] Symmetric encryption
		\item[\textbf{(C)}] A hash function
		\item[\textbf{(D)}] A digital signature}
	{B}
	{\item Symmetric algorithms are fast for bulk data.
		\item Asymmetric methods are often used to exchange or protect the symmetric keys.
		\item Hashes are one-way and cannot be decrypted.}
	{Symmetric for speed, asymmetric for key exchange.}
	
	\hq{91}{Checking File Integrity}
	{How can a recipient confirm that a downloaded file was not altered?}
	{\item[\textbf{(A)}] Check the file size only
		\item[\textbf{(B)}] Open it and look for errors
		\item[\textbf{(C)}] Rename it and scan again
		\item[\textbf{(D)}] Compare its cryptographic hash value with the value published by the source}
	{D}
	{\item Even a tiny change produces a different hash.
		\item The reference hash must come from a trusted channel.
		\item Size and appearance are unreliable checks.}
	{A matching hash means unchanged.}
	
	\hq{92}{What a Signature Provides}
	{Which combination does a digital signature provide?}
	{\item[\textbf{(A)}] Confidentiality only
		\item[\textbf{(B)}] Availability and performance
		\item[\textbf{(C)}] Anonymity of the sender
		\item[\textbf{(D)}] Authentication of the sender, integrity of the message, and non-repudiation}
	{D}
	{\item The signature is created with the sender's private key.
		\item Any change to the message invalidates it.
		\item It does not by itself encrypt the content.}
	{A signature proves who sent it and that it did not change.}
	
	\hq{93}{Certificate Expiry Outage}
	{A critical website went down because a certificate expired unnoticed. What is the \textbf{BEST} preventive measure?}
	{\item[\textbf{(A)}] Maintain a certificate inventory with expiry monitoring and renewal ownership
		\item[\textbf{(B)}] Buy certificates that never expire
		\item[\textbf{(C)}] Ignore certificates for internal systems
		\item[\textbf{(D)}] Wait for users to report problems}
	{A}
	{\item Expiry is predictable and preventable.
		\item Ownership ensures renewals are done.
		\item Monitoring warns before service is affected.}
	{Certificates expire on schedule; plan for it.}
	
	\hq{94}{Key Storage}
	{Encryption keys are stored on the same server as the data they protect. What is the \textbf{BEST} recommendation?}
	{\item[\textbf{(A)}] Store keys separately in a hardened key management system or hardware security module
		\item[\textbf{(B)}] Leave them because encryption is in place
		\item[\textbf{(C)}] Email the keys to the administrators
		\item[\textbf{(D)}] Store keys in a shared spreadsheet}
	{A}
	{\item If an attacker gets both, encryption gives no protection.
		\item Separate, protected storage limits exposure.
		\item Key access must be restricted and logged.}
	{Do not lock the door and tape the key to it.}
	
	\hq{95}{Effective Awareness Training}
	{Which approach to security awareness training is \textbf{MOST} effective?}
	{\item[\textbf{(A)}] One annual video for everyone
		\item[\textbf{(B)}] Regular, role-based training with measured outcomes such as phishing simulations
		\item[\textbf{(C)}] Training only for the IT department
		\item[\textbf{(D)}] Training only after a breach}
	{B}
	{\item Different roles face different risks.
		\item Repetition reinforces behavior.
		\item Metrics show whether behavior is improving.}
	{Train often, tailor it, and measure it.}
	
	\hq{96}{Data Minimization}
	{Marketing wants to collect many personal data fields from customers just in case they are useful later. What should the risk practitioner recommend?}
	{\item[\textbf{(A)}] Collect everything while it is possible
		\item[\textbf{(B)}] Collect more only from new customers
		\item[\textbf{(C)}] Collect only the data necessary for a defined, lawful purpose
		\item[\textbf{(D)}] Let marketing decide without review}
	{C}
	{\item Less data means less exposure and less compliance burden.
		\item Purpose should be defined before collection.
		\item Excess data increases breach impact.}
	{Collect what you need, not what you can.}
	
	\hq{97}{Privacy Impact Assessment}
	{A new application will process personal data. What should be done \textbf{BEFORE} go-live?}
	{\item[\textbf{(A)}] Wait for the first customer complaint
		\item[\textbf{(B)}] Perform a privacy impact assessment and address the findings
		\item[\textbf{(C)}] Delay all privacy work until after launch
		\item[\textbf{(D)}] Assume the security review covers privacy}
	{B}
	{\item A PIA identifies privacy risks early.
		\item Findings may change design or data handling.
		\item Security and privacy overlap but are not the same.}
	{Assess privacy before launch, not after complaints.}
	
	\hq{98}{Deletion Request}
	{A customer asks the company to delete all of their personal data. What is the \textbf{BEST} approach?}
	{\item[\textbf{(A)}] Verify identity, check legal retention duties, and respond within the required timeframe
		\item[\textbf{(B)}] Delete without any verification
		\item[\textbf{(C)}] Ignore the request
		\item[\textbf{(D)}] Ask the customer to call next year}
	{A}
	{\item Requests must come from the verified data subject.
		\item Some data may need to be retained by law.
		\item Regulations define response deadlines.}
	{Honor privacy rights through a defined, verified process.}
	
	\hq{99}{Security vs Privacy}
	{Which statement \textbf{BEST} distinguishes data security from data privacy?}
	{\item[\textbf{(A)}] Security protects data from unauthorized access, while privacy governs appropriate collection, use, and sharing of personal data
		\item[\textbf{(B)}] They are the same thing
		\item[\textbf{(C)}] Privacy applies only to encrypted data
		\item[\textbf{(D)}] Security applies only to paper records}
	{A}
	{\item A system can be secure yet misuse personal data.
		\item Privacy requires lawful and fair processing.
		\item Both are needed, and they work together.}
	{Secure is not the same as private.}
	
	\hq{100}{Cross-Border Data Transfer}
	{The company plans to move customer personal data to a data center in another country. What should be done \textbf{FIRST}?}
	{\item[\textbf{(A)}] Move the data first and ask legal later
		\item[\textbf{(B)}] Choose the cheapest country
		\item[\textbf{(C)}] Assume all countries have identical rules
		\item[\textbf{(D)}] Assess applicable laws and legal transfer mechanisms before moving the data}
	{D}
	{\item Data transfer rules vary by jurisdiction.
		\item Non-compliance can lead to fines and forced stoppage.
		\item Contracts and safeguards may be required.}
	{Know the law of the place your data lives.}
	
	% ============================================================
	\section{Conclusion}
	% ============================================================
	
	\begin{frame}{Exam Mindset Summary}
		\begin{block}{Remember These Core Principles}
			\begin{itemize}
				\item \textbf{Owners Decide:} The risk owner accepts, mitigates, transfers, or avoids risk; practitioners advise.
				\item \textbf{Fit the Response:} Choose responses that bring risk within appetite at reasonable cost.
				\item \textbf{Controls Serve Risks:} Select, implement, and test controls based on identified risks.
				\item \textbf{Verify, Do Not Assume:} Evidence of effectiveness matters more than assertions or tool purchases.
				\item \textbf{Monitor Continuously:} Use well-chosen indicators tied to appetite to trigger timely action.
				\item \textbf{Understand Before Acting:} Know the architecture, data, and business context before choosing technical fixes.
				\item \textbf{Accountability Cannot be Outsourced:} Vendors and cloud providers share duties, but accountability remains with the organization.
			\end{itemize}
		\end{block}
	\end{frame}
	
	\begin{frame}{Thank You}
		\begin{center}
			{\Huge \textbf{Good Luck on Your CRISC Exam!}}
			
			\vspace{1cm}
			
			\normalsize
			Real exam questions test your ability to apply these concepts in practical, real-world scenarios.
		\end{center}
	\end{frame}
	
\end{document}